\section{Introduction}
\label{sec:intro}

Sleep is one of very few \emph{modifiable} candidate risk factors for Alzheimer's disease, and
the glymphatic hypothesis gives it a concrete mechanism. Slow waves drive macroscopic
cerebrospinal-fluid pulsation in humans \citep{fultz2019}, that flow clears interstitial
amyloid-beta in mice \citep{xie2013}, and so chronically reduced slow-wave sleep should permit
amyloid to accumulate. The chain is repeated in reviews, in press coverage, and --- most
consequentially --- in the design of acoustic-stimulation and behavioural sleep trials that
treat N3 sleep duration as the actionable target. If the human evidence for the \emph{stage}
metric is weaker than the narrative implies, those trials are being powered and targeted on an
effect that may not exist at the claimed size.

\para{The hypothesis is a conjunction of three claims.} As usually stated, reduced slow-wave
sleep duration shows a dose-dependent association with higher amyloid burden in cognitively
normal adults, each 10\% reduction in N3 percentage is associated with a measurable increase in
amyloid PET SUVR, and this relationship is stronger than the associations with total sleep time
or sleep efficiency alone. Clause (i) is an association claim, clause (ii) a dose claim in
physical units, and clause (iii) a comparative claim. Each is separately testable, and the
hypothesis is supported only if all three hold.

\para{None of the three has ever been tested meta-analytically.} From 54 full-text papers we
find four gaps, and every one is load-bearing. First, no meta-analysis has pooled a
slow-wave-specific exposure against amyloid PET: the largest and most recent synthesis
\citep{chen2025} pooled only sleep \emph{quality} and \emph{duration} across 30 studies and
14{,}997 subjects, while \citet{cui2022} and \citet{bubu2020} covered obstructive sleep apnea.
Second, ``slow-wave sleep'' silently names two different measurements. In the one dataset that
reports both in the same 32 participants, spectral $<$1\,Hz slow-wave activity predicted amyloid
accumulation at $r = -0.52$ while N3 \emph{stage} duration was null at $r = -0.07$
\citep{winer2020}; any estimate that merges them is uninterpretable. Third, cohort
non-independence is unhandled: \citet{mander2015}, \citet{winer2019} and \citet{winer2020} are
all the Berkeley Aging Cohort Study with overlapping participants, contributing 13 of our 49
extracted effects. Fourth, and least appreciated, differential measurement reliability has
never been tested as a competing explanation for why one sleep metric correlates more strongly
than another.

\para{Our approach.} We assemble 49 hand-extracted effect sizes from 17 studies across 12
independent cohort families, each carrying a verified DOI/PMID, a verbatim provenance quote, and
a sign convention in which \emph{positive means worse sleep is associated with more amyloid}.
We pool the 40 amyloid-PET effects on the Fisher-$z$ scale in four preregistered strata ---
\stage, \spectral, \tst and \se{} --- aggregating to one synthetic effect per cohort family
before fitting REML random effects with Knapp--Hartung--Sidik--Jonkman standard errors. We then
test the comparative claim as a \emph{within-study paired difference} using only cohorts that
report both exposures in the same participants, build a measurement-reliability null model from
empirical single-night ICCs, and convert the pooled estimate into SUVR and Centiloid units
(\figref{fig:forest} through \figref{fig:robust}).

\para{What we find.} The exposure the hypothesis names is the weakest of the four:
$r = +0.026$ (95\% CI $-0.153$ to $+0.203$, $p = 0.68$) with $I^2 = 0\%$, against $r = +0.163$
for total sleep time and $r = +0.090$ for sleep efficiency. Equivalence testing rejects a true
correlation as large as $r = 0.20$ ($p_{\mathrm{TOST}} = 0.026$). The dose statement converts to
$\Delta\mathrm{SUVR} = +0.005$ per 10 \pp of N3, $1.8\%$ of the amyloid-negative/positive
separation. Measurement error cannot rescue the comparative claim: single-night polysomnography
measures N3\% far more reliably (ICC $0.775$) than total sleep time (ICC $0.225$), so identical
true effects would still make N3 look $1.86\times$ stronger --- the observed ratio is
$0.16\times$, about eleven times below the measurement-only prediction. Finally, the positive
spectral literature is one cohort family: pooled $r = +0.112$ across three families
($I^2 = 77\%$) drops to $r = -0.025$ when Berkeley is removed.

\para{Contributions.}
\begin{itemize}[leftmargin=*,itemsep=1pt,topsep=2pt]
    \item We provide the first meta-analytic estimate of the N3-stage/amyloid-PET association,
    kept strictly separate from spectral slow-wave activity and with variance handled at the
    level of the cohort family rather than the publication.
    \item We test the specificity claim as a within-study paired contrast, using only cohorts
    that measure both exposures in the same participants, which removes between-cohort
    confounding from the exact comparison the hypothesis makes.
    \item We introduce a quantitative measurement-reliability null model from empirical
    single-night ICCs, turning ``metric A correlates more strongly than metric B'' into a
    testable excess over what measurement precision alone predicts.
    \item We estimate the hypothesis's own dose statement in SUVR and Centiloid units, and show
    with population N3 distributions that the stated 10-percentage-point contrast is not
    available to 78\% of the target population.
\end{itemize}

\Secref{sec:related} positions this work against the mechanistic, observational and
meta-analytic literature; \secref{sec:method} gives the harmonisation, dependency handling and
six preregistered analyses; \secref{sec:results} reports them; \secref{sec:discussion} states
the verdict on each clause, what the reliability model adds, and the limits of a study-level
design.
